%% main.tex — transmon-benchmark.3 (anvil:pub revision, iteration 3)
%%
%% Revised 2026-07-14 by pub-revise from transmon-benchmark.2 + the
%% transmon-benchmark.2.litsearch critic sibling, against the HEAVILY
%% UPDATED papers/transmon-benchmark/BRIEF.md (thesis reframe to
%% tensor-compiler viability; framing gate resolved to Branch A; GPU
%% scaling data final; new related-work axis; authors resolved).
%% No transmon-benchmark.2.review sibling exists — this revision is
%% operator-directed via the BRIEF delta; see changelog.md.
%%
%% Every agreement number is quoted VERBATIM from the committed
%% benchmarks/transmon_eigen/results.toml. CPU-cell numbers trace to
%% benchmarks/transmon_bench_cpu/results.toml. GPU driven-solve scaling
%% numbers trace to benchmarks/gpu_driven_scaling/results.toml (merged,
%% PR #516). Gauge/projection-arc and spurious-mode-characterization
%% numbers trace to the [spurious_mode] block of
%% benchmarks/transmon_eigen/results.toml (PRs #508/#513/#515). Quantum
%% parameters trace to benchmarks/transmon_quantum/results.toml (PR #511).
%% See changelog.md in this version dir for the note -> change map.

\documentclass{anvil-paper}

%% --- URL line-breaking: the two AWS-blog references carry very long
%% hyphenated URLs; allow breaks at hyphens and add stretch so the
%% bibliography stays inside the measure (render-gate overfull check).
\expandafter\def\expandafter\UrlBreaks\expandafter{\UrlBreaks\do\-}
\Urlmuskip=0mu plus 1mu

%% --- Draft-phase helper macros -------------------------------------------
%% Figure placeholder: uses the rendered figure if pub-figures has produced
%% it, otherwise a framed placeholder naming the source script.
\newcommand{\anvilfig}[3][0.85]{%
  \IfFileExists{figures/#2}%
    {\includegraphics[width=#1\linewidth]{figures/#2}}%
    {\fbox{\parbox{#1\linewidth}{\centering\vspace{5em}%
      \texttt{figures/#2}\\[0.5em]%
      {\small placeholder --- render via \texttt{figures/src/#3}}%
      \vspace{5em}}}}}

\title{Tensor-compiler--based finite-element electromagnetics:
cross-validating a Burn-native $H(\mathrm{curl})$ solver against Palace on
a superconducting transmon}
% TODO(operator): final title wording is operator-approved per BRIEF.md
% "Framing consequences"; this is the BRIEF's suggested Branch-A wording.
\author{Robb Walters\thanks{These authors contributed equally.}
  \and Crutcher Dunnavant\thanks{These authors contributed equally.}}
% TODO(operator): affiliations (author order and equal-contribution
% footnote are operator-resolved per BRIEF.md, 2026-07-14).
\date{\today}

\begin{document}
\maketitle

\begin{abstract}
The machine-learning ecosystem has produced general-purpose tensor-compiler
stacks --- batched kernels, JIT compilation, transparent multi-backend
portability --- representing one of the largest sustained engineering
investments in numerical computing. We ask whether that investment can be
inherited by classical mesh-based electromagnetics, and answer with a
measured, honestly bounded yes. GEODE-FEM is a full-wave 3D
$H(\mathrm{curl})$ finite-element solver built in Rust on the Burn tensor
framework: element assembly is batched $[n_{\mathrm{elem}},6,6]$ tensor
algebra, the matrix-free operator is a gather $\to$ batched-matmul $\to$
scatter-add pipeline, and the driven Krylov iteration runs on-device with an
$O(1)$-scalar per-iteration synchronization budget, retargeting CPU and CUDA
from one codebase. The claim is held to a cross-validation evidence
standard: on a real transmon-plus-readout-resonator geometry meshed once and
consumed identically by both solvers, GEODE-FEM reproduces Palace --- the
MFEM-based reference solver for superconducting-qubit design --- to within
$0.033\%$ across all six eigenmodes (worst case $0.032\%$; junction LC mode
$0.001\%$), with exact analytic tripwires guarding against coincidental
agreement. Performance is reported as measured, wins and losses alike: the
serial direct-factorization eigensolve is ${\sim}4\times$ more efficient per
core than Palace's MPI-distributed solve at 133k DOFs (Palace wins
$1.7\times$ at the whole-node level), while the GPU driven-solve scaling
cell is an honest negative --- at $\leq 25{,}695$ edges the CUDA-f32
matrix-free path loses to every CPU configuration (assembled-CSR COCG is
$44\times$ faster at the top size), reaching parity only against its own
algorithm on CPU at ${\sim}26$k edges, with f32 accuracy flooring near
$10^{-3}$. Honest physics is the spine throughout: the model contains no
${\sim}4$\,GHz qubit mode \emph{by construction}, and a disclosed spurious
mode is diagnosed through a three-step gauge/projection arc --- tree--cotree
elimination (spectrum-shifting, $1.64\%$), bulk divergence-free projection
(cavity-preserving but junction-deflating), and a port-aware rank-1
re-admission that lands all six modes at $\leq 0.029\%$ --- and characterized
as a port-surface artifact. Every number traces to a committed artifact.
\end{abstract}

\section{Introduction}
\label{sec:intro}

The tensor-compiler stacks built for machine learning --- PyTorch, JAX,
and, in Rust, Burn \citep{burn} --- embody an enormous, continuously
maintained engineering investment: JIT kernel generation, batched dense
linear algebra, memory planners, and single-source portability across CPU,
CUDA, and WebGPU backends. That infrastructure was built for neural
networks, but nothing in it is specific to them. This paper asks a
systems question: \emph{can a production-grade computational
electromagnetics solver be built on a general-purpose ML tensor-compiler
stack} --- inheriting its kernels, its JIT, and its backend portability
--- \emph{rather than on the domain-specific FEM infrastructure the field
has built for itself?} And it insists on an unforgiving evidence standard
for the answer: same-mesh, same-physics cross-validation against the
reference solver, at production accuracy, with every number traced to a
committed artifact.

The mapping is natural once stated. Finite-element assembly is
embarrassingly element-parallel dense algebra: for first-order N\'ed\'elec
elements, the curl--curl and mass contributions are per-tetrahedron
$6\times 6$ matrices, which a tensor framework expresses as one batched
$[n_{\mathrm{elem}},6,6]$ kernel. Matrix-free operator application is a
gather (global edge DOFs to element-local vectors), a batched matrix
multiply, and a scatter-add --- exactly the primitives ML frameworks
optimize. A Krylov iteration is a short loop of such applies plus
reductions, which can stay resident on an accelerator with a constant
number of scalar host--device synchronizations per iteration. GEODE-FEM
(``geode-fem'') realizes this mapping in Rust on Burn
(Section~\ref{sec:architecture}), with honest constraints stated up front:
the CUDA backend is f32-only today, and the sparse direct factorization at
the heart of the eigensolve remains a classical CPU component outside the
tensor-compiler story.

The validation vehicle is deliberately demanding. Palace \citep{palace},
an MFEM-based \citep{mfem,andrej2024high} parallel solver from the AWS
Center for Quantum Computing, is the open-source solver of record for
superconducting-qubit electromagnetic design, typically driven by geometry
from DeviceLayout.jl \citep{devicelayout,devicelayout-blog}. Mode
frequencies and participation ratios from such simulations feed directly
into device Hamiltonians \citep{koch2007,nigg2012,minev2021}, so the
accuracy bar is production-grade, not demonstration-grade. Notably, this
influential stack has never been the subject of a peer-reviewed or
preprint publication --- Palace exists as a repository, a blog post
\citep{palace-blog}, and a workshop talk; the DeviceLayout.jl transmon
workflow as a blog post and a conference talk
\citep{devicelayout-blog,devicelayout-juliacon} --- and no independent,
quantitative, third-party validation of it exists in the literature. This
paper therefore does double duty: the cross-validation that anchors the
architecture claim is also the first independent validation of the
reference stack itself. The comparison is same-mesh, same-physics,
same-junction-model --- the geometry is meshed \emph{once} and both
solvers consume the identical mesh, first-order discretization, and lumped
reactive-shunt junction --- so agreement claims are about formulation and
solver correctness, not meshing luck.

Our contributions are:
\begin{itemize}
  \item \textbf{Evidence that a general-purpose ML tensor stack can carry
    production-accuracy computational electromagnetics.} To our knowledge,
    geode-fem is the first full-wave 3D $H(\mathrm{curl})$/N\'ed\'elec FEM
    solver built on a general-purpose ML tensor-compiler stack and
    validated at production accuracy against the reference solver: on the
    identical mesh at matched order it reproduces all six of Palace's
    eigenmodes to $\leq 0.033\%$ (worst case $0.032\%$), with the junction
    LC mode agreeing to $0.001\%$ (Section~\ref{sec:agreement}). The same
    result is the first independent validation of the
    Palace/DeviceLayout.jl transmon stack.
  \item \textbf{An architecture for FEM as a tensor-compiler workload.}
    Batched element-local assembly, a matrix-free
    gather$\to$batched-matmul$\to$scatter-add operator, and an on-device
    Krylov loop with an $O(1)$-scalar synchronization budget, retargeting
    CPU/CUDA from one codebase --- with the current boundaries (f32-only
    CUDA; factorization-bound eigensolve) stated as constraints, not
    footnotes (Section~\ref{sec:architecture}).
  \item \textbf{Honest performance accounting on both axes.} On CPU, the
    serial direct-factorization eigensolve is ${\sim}4\times$ more
    efficient per core than Palace's distributed Krylov solve at 133k
    DOFs, while MPI parallelism wins $1.7\times$ at the whole-box level
    (Section~\ref{sec:perf-cpu}). On GPU, a same-host four-configuration
    scaling study is reported as the honest negative it is: the CUDA-f32
    matrix-free driven solve loses to every CPU configuration at every
    measured size, with a single directional positive --- monotonic
    convergence to parity against the same algorithm on CPU at ${\sim}26$k
    edges (Section~\ref{sec:perf-gpu}).
  \item \textbf{A reproducible same-mesh cross-solver methodology.}
    Oracle-first benchmarking --- exact analytic tripwires, inverse tests
    that must fail, participation-based mode identification, and a fully
    scripted fixture-to-comparison pipeline
    (Sections~\ref{sec:methodology} and~\ref{sec:repro}).
  \item \textbf{Honest-negative reporting as a first-class result.} The
    model's lack of a ${\sim}4$\,GHz qubit mode is explained and confirmed
    by both solvers (Section~\ref{sec:absent-mode}); a spurious mode
    leaked by geode-fem's un-projected Lanczos is disclosed,
    \emph{diagnosed through a measured three-step gauge/projection arc},
    characterized as a port-surface artifact, and resolved by a port-aware
    projection that retains all six physical modes
    (Section~\ref{sec:spurious-mode}).
\end{itemize}

Section~\ref{sec:related} positions this work along both axes ---
FEM-on-tensor-stacks and qubit-EM validation; Section~\ref{sec:architecture}
presents the architecture; Sections~\ref{sec:problem}--\ref{sec:solvers}
define the benchmark problem and the solver contrast; the remaining
sections present methodology, results, and reproducibility details.

\section{Related Work}
\label{sec:related}

\paragraph{FEM by code generation --- and the libCEED foil.}
Generating high-performance FEM kernels from high-level descriptions is a
two-decade tradition: the FFC form compiler \citep{kirby2006compiler}, the
UFL weak-form language \citep{alnaes2014unified}, Firedrake's
composition-of-abstractions approach \citep{rathgeber2017firedrake}, and
the structure-preserving TSFC compiler \citep{homolya2018tsfc} all compile
variational forms to optimized element kernels. The modern GPU expression
of this tradition is libCEED \citep{brown2021libceed} --- matrix-free,
element-based operator evaluation --- together with MFEM's partial-assembly
GPU path \citep{andrej2024high}. The positioning matters for this paper:
\emph{Palace itself runs on that stack} --- MFEM with libCEED-class
backends, a domain-specific element-kernel JIT built by and for the FEM
community. geode-fem's bet is categorically different: outsource kernel
generation, scheduling, and backend portability to a \emph{general-purpose
ML tensor stack} whose engineering is amortized across the entire ML
ecosystem, and phrase assembly and matrix-free application as batched
tensor algebra. Form compilers generate FEM kernels; the tensor-compiler
approach reuses someone else's compiler and makes FEM look like its native
workload. The comparison in this paper is therefore not
``hand-written versus generated'' but ``domain-specific generator versus
general-purpose stack'' --- with the same physics on both sides.

\paragraph{General-purpose tensor compilers and FEM on ML frameworks.}
The compiler lineage geode-fem's substrate inherits runs through Halide's
algorithm/schedule decoupling \citep{ragankelley2013halide} and TVM's
end-to-end tensor-program optimization \citep{chen2018tvm}; we found no
peer-reviewed EM or FEM solver built directly on that lineage --- its
scientific-computing uptake has come via the ML frameworks. There,
JAX-FEM \citep{xue2023jax} is the canonical example (nodal-Lagrange
continuum mechanics on JAX, GPU via XLA, differentiable), with PyTorch-FEA
\citep{liang2023pytorch} the biomechanics analog. Closest to this work is
the concurrent TensorMesh/TensorGalerkin framework \citep{wen2026learning}
with its companion sparse-linear-algebra library torch-sla
\citep{chi2026torch}: general differentiable Galerkin assembly as batched
tensor operations on PyTorch, whose demonstration suite includes
magnetostatics via a stabilized \emph{nodal} curl--curl formulation. Its
existence independently supports the architectural thesis --- FEM assembly
as batched tensor algebra on an ML framework is now a genus, not a one-off
--- while sharpening what remains distinct here: TensorMesh is nodal
rather than $H(\mathrm{curl})$-conforming, its electromagnetic examples
are magnetostatic rather than full-wave generalized eigenproblems with
lumped-element loading, and it reports no cross-validation against a
reference EM solver at production accuracy. To our knowledge, no
\emph{full-wave 3D $H(\mathrm{curl})$/N\'ed\'elec} FEM solver has been
built on a general-purpose ML tensor stack \emph{and validated at
production accuracy against the reference solver} before this work.

\paragraph{Differentiable and GPU electromagnetics.}
The differentiable-EM literature reaches Maxwell through structured grids:
autograd FDFD/FDTD in ceviche \citep{hughes2019forward}, JAX-based inverse
design under foundry constraints \citep{schubert2022inverse}, the Meep
adjoint pipeline (autograd/JAX around a classical FDTD core)
\citep{hammond2022high}, the JAX-native FDTDX framework
\citep{mahlau2026fdtdx}, and differentiable RCWA in Meent
\citep{kim2024meent}. These codes get Maxwell onto ML frameworks by
accepting rectilinear grids and low-order stencils; the transmon problem
requires body-fitted tetrahedra, anisotropic sapphire, and
$H(\mathrm{curl})$ conformity --- exactly the regime that literature has
not reached. Adjacent but distinct genera: purpose-built differentiable
kernel DSLs (DiffTaichi \citep{hu2019difftaichi}; $\Phi_{\mathrm{Flow}}$
\citep{holl2020learning}), where the language is built \emph{for}
simulation rather than reused from ML; neural-ansatz Maxwell eigensolvers
such as FieldTNN \citep{jiang2024fieldtnn}, where a network \emph{replaces}
the discretization rather than executing it; and the hand-written-CUDA GPU
FEM tradition for electromagnetics \citep{huantingmeng2014gpu}, which
targets the same hardware at the cost of bespoke kernel engineering ---
the specialization-for-portability trade the tensor-compiler approach
inverts.

\paragraph{Electromagnetic simulation workflows for superconducting qubits.}
The closest prior work on the validation axis consists of two recent
Palace-based workflow papers. \citet{sommers2025open} present SQDMetal,
an open workflow coupling Qiskit Metal, Gmsh, and Palace, and cross-check
Palace against the commercial solvers COMSOL and Ansys HFSS on qubit and
resonator geometries, reporting agreement at roughly the $0.02$--$0.13\%$
level. \citet{ye2025electromagnetic} build a Palace-centered
layout-to-Hamiltonian workflow and benchmark it against cryogenic
measurement, with resonator frequencies reproduced to within $0.3\%$. Both
efforts treat Palace as a component to be \emph{wrapped and trusted}:
validation runs against commercial tools or experiment, through
independently constructed models and meshes, so residuals mix solver error
with meshing and model-entry differences. Neither is a same-mesh,
same-formulation comparison against an independently \emph{implemented}
solver. Our study eliminates meshing and model entry as confounders ---
identical MSH fixture, identical lumped-shunt junction, first-order
N\'ed\'elec elements on both sides --- so residuals isolate formulation
and solver correctness; to our knowledge it is the first independent
third-party validation of Palace itself. The broader open ecosystem for
qubit electromagnetic design includes Qiskit Metal and its quasi-lumped
analysis methods \citep{minev2021circuit} and the SQuADDS validated design
database \citep{shanto2024squadds}; the physics framing for extracting
qubit parameters from field modes rests on the transmon \citep{koch2007},
black-box quantization \citep{nigg2012}, and energy-participation-ratio
\citep{minev2021} literature.

\paragraph{Verification, validation, and cross-solver benchmarking.}
Code-to-code comparison on precisely specified benchmark problems is an
established genre in computational electromagnetics: the TEAM (Testing
Electromagnetic Analysis Methods) workshop tradition has run such
comparisons since 1986.\footnote{The TEAM founding documents are workshop
reports and proceedings without resolvable DOIs; we describe the tradition
in prose rather than cite an arbitrary secondary source.} In the V\&V
vocabulary of \citet{roache1997quantification} and
\citet{oberkampf2010verification}, the present study is
verification-side evidence: our analytic tripwires are code-verification
instruments in exactly Roache's sense, and the same-mesh protocol is a
code-to-code comparison in Oberkampf and Roy's taxonomy. Multi-solver
comparisons on shared electromagnetic problems have precedent in
nano-optics, where \citet{hoffmann2009comparison} compared FEM, FDTD, and
FIT codes on plasmonic antennas. We import that discipline to the
superconducting-qubit domain and add two elements the genre often lacks:
exact analytic tripwires that must move in known ways, and honest-negative
reporting of what the model does \emph{not} contain.

\paragraph{Open-source solver ecosystems and numerics.}
Mature open electromagnetics tools are normally accompanied by a citable
methods paper: GetDP \citep{dular1998general}, DOLFIN/FEniCS
\citep{logg2010dolfin}, Gmsh \citep{gmsh}, and scikit-rf
\citep{arsenovic2022scikit} all have one; so does Palace's substrate
library MFEM \citep{mfem,andrej2024high}. Palace and DeviceLayout.jl are
outliers --- the \emph{library} is published, the solver product built on
it is not --- which is precisely the gap this benchmark narrows from the
outside. On the numerics side, geode-fem's eigenpath is a shift-invert
Lanczos \citep{arpack} over a real symmetric pencil, factorized with the
faer sparse LU \citep{kazdadi2026faer}, with first-order N\'ed\'elec
elements \citep{nedelec1980}; Palace's eigenpath is a SLEPc-class Krylov
method \citep{slepc} with divergence-free projection and auxiliary-space
(AMS) preconditioning
\citep{hiptmair-xu,tzaniovkolevpanayotsvassilevski2018parallel}. The
gauge-fixing literature relevant to Section~\ref{sec:spurious-mode} is the
tree--cotree tradition \citep{albanese1988solution,manges1995generalized}.

\section{GEODE-FEM: Finite Elements as a Tensor-Compiler Workload}
\label{sec:architecture}

geode-fem is written in Rust on the Burn tensor framework \citep{burn},
whose cubecl JIT compiles tensor kernels for CPU, CUDA, and WebGPU
backends from one source. The design principle is to phrase every
throughput-critical FEM stage as the batched dense algebra the stack
already optimizes, and to be explicit about which stages do not (yet) fit.

\paragraph{Batched element-local assembly.}
For first-order N\'ed\'elec tetrahedra the element curl--curl and mass
contributions are $6\times 6$ matrices depending on per-element geometry
factors and material tensors (including the rotated anisotropic sapphire
permittivity of Section~\ref{sec:geometry}). geode-fem computes them for
all elements at once as batched $[n_{\mathrm{elem}},6,6]$ tensor
operations --- one kernel launch per assembly stage rather than an
element loop --- then either scatters into a global sparse matrix (the
assembled path used by the f64 CPU eigensolve) or retains the batch for
matrix-free application.

\paragraph{Matrix-free operator application.}
The matrix-free apply of $(K + \mathrm{i}\omega\text{-terms})x$ is a
three-stage tensor pipeline: \emph{gather} global edge coefficients into
element-local $[n_{\mathrm{elem}},6]$ vectors, \emph{batched matmul}
against the stored $[n_{\mathrm{elem}},6,6]$ element operators, and
\emph{scatter-add} back to global DOFs. Each stage is a native
tensor-framework primitive; the backend (ndarray on CPU, CUDA on GPU)
is a type parameter, not a code path.

\paragraph{On-device Krylov with an $O(1)$-scalar sync budget.}
The driven (deterministic frequency-sweep) solve runs a conjugate
orthogonal conjugate gradient (COCG) iteration entirely on the device:
matrix-free applies, axpy updates, and inner-product reductions all
execute where the data lives, and only a constant number of scalars per
iteration (convergence and recurrence coefficients) cross the host--device
boundary. This is the architectural bet that device-resident Krylov
amortizes once problems are large enough to fill an accelerator ---
Section~\ref{sec:perf-gpu} measures, honestly, how far below that
threshold the current benchmark sizes sit.

\paragraph{Honest constraints.}
Two boundaries of the tensor-compiler story are stated as constraints
rather than buried. First, the CUDA backend (burn-cuda 0.21) is f32-only
today --- the underlying cubecl runtime currently disables
f64\footnote{Upstream constraint of the burn/cubecl stack at version 0.21
\citep{burn}; tracking-issue URL to be added at revision.
TODO(operator).} --- so every GPU result in this paper is
mixed-precision-qualified and reported with explicit accuracy deltas
against the f64 CPU reference. Second, the eigensolve that produces the
headline comparison is factorization-bound: its inner solves are faer
sparse-LU factorizations \citep{kazdadi2026faer}, a classical serial CPU
component. The tensor-compiler story currently covers assembly and driven
solves, \emph{not} sparse direct factorization; the eigensolver
gauge/projection work of Section~\ref{sec:spurious-mode} (repository
issues \#502/\#503 and follow-ons) is a prerequisite for a future
projected, iterative --- and therefore GPU-viable --- eigenpath.

\section{Benchmark Problem: Geometry, Mesh, and Junction Model}
\label{sec:problem}

\subsection{Geometry and mesh provenance}
\label{sec:geometry}

The device is the \texttt{SingleTransmon} example shipped with
DeviceLayout.jl v1.15.0 \citep{devicelayout}: a transmon with readout
resonator on a sapphire substrate, enclosed in a vacuum box. The sapphire is
anisotropic, with permittivity tensor
$\varepsilon = R\,\mathrm{diag}(9.3,\, 9.3,\, 11.5)\,R^{\mathsf{T}}$
rotated approximately $36.87^\circ$ in-plane. The DeviceLayout.jl workflow
emits seven named physical groups (metal, substrate, vacuum, junction
lumped-element surface, two readout ports, exterior boundary) and drives
Gmsh \citep{gmsh} to produce a tetrahedral mesh, exported as an MSH~4.1
fixture and pinned by SHA-256 (Section~\ref{sec:repro}).

The mesh has 22{,}684 nodes and 133{,}314 tetrahedra, yielding 156{,}863
N\'ed\'elec edge unknowns, of which 133{,}108 interior degrees of freedom
remain after eliminating perfect-electric-conductor (PEC) boundary edges.
PEC is applied on all metal surfaces and the exterior box; the two readout
ports are left open (lossless) in this first version, so no resistive
element enters and the eigenproblem pencil stays real symmetric.
Figure~\ref{fig:geometry} shows the geometry with physical groups
color-coded.

\begin{figure}[t]
  \centering
  \anvilfig{fig1-geometry.pdf}{fig1\_geometry.py}
  \caption{Benchmark geometry: top-down projection of the committed
  surface mesh (the SHA-256-pinned MSH~4.1 fixture) for the
  DeviceLayout.jl v1.15.0 \texttt{SingleTransmon} transmon and readout
  resonator. Surface physical groups are color-coded --- the metal ground
  plane (with the coplanar-waveguide feedline, meander resonator, and
  transmon slot as etched gaps), the four-triangle junction
  \texttt{lumped\_element} surface (enlarged in the inset, actual
  triangulation shown), the two readout ports, and the exterior-boundary
  outline; the \texttt{substrate} and \texttt{vacuum} volume groups fill
  the enclosing box. Both solvers consume this identical mesh (22{,}684
  nodes; 133{,}314 tetrahedra).}
  \label{fig:geometry}
\end{figure}

\subsection{The junction model: a lumped reactive shunt}
\label{sec:junction}

The Josephson junction is modeled identically in both solvers as a lumped
reactive shunt --- a parallel $L \parallel C$ element with
$L = 14.860$\,nH and $C = 5.5$\,fF --- attached to the four-triangle
\texttt{lumped\_element} surface group $\Gamma$. Following the uniform
lumped-port construction, with port length $\ell$ and width $w$, the shunt
contributes
\begin{equation}
  K_{\mathrm{port}} = \frac{\ell}{w\,L}\, S_\Gamma
  \qquad\text{and}\qquad
  M_{\mathrm{port}} = \frac{C\,\ell}{w}\, S_\Gamma ,
  \label{eq:port}
\end{equation}
where $S_\Gamma$ is the port-surface matrix assembled on $\Gamma$. The
inductive term is frequency-independent and adds to the stiffness matrix;
the capacitive term adds to the mass matrix. The eigenproblem is then the
real symmetric generalized pencil
\begin{equation}
  \left(K + K_{\mathrm{port}}\right) x
  \;=\; \omega^2 \left(M + M_{\mathrm{port}}\right) x .
  \label{eq:pencil}
\end{equation}
In Palace, the same physics is entered as a reactive \texttt{LumpedPort}
($L \parallel C$, no resistance) on the same surface group; the readout
ports are configured open (no $50\,\Omega$ termination), so Palace's pencil
is likewise real. This one-to-one correspondence of the junction treatment
is what makes the comparison same-physics as well as same-mesh.

\section{Solver Architectures Under Comparison}
\label{sec:solvers}

The two solvers share the mathematical problem
(Eq.~\ref{eq:pencil}) and nothing else: no code, no language, no linear
algebra, no eigensolver strategy --- and, per Section~\ref{sec:related},
they sit on opposite sides of the domain-specific-versus-general-purpose
compiler divide. Table~\ref{tab:architectures} summarizes the contrast.

\begin{table}[t]
  \centering
  \caption{Solver architectures compared. The two implementations share the
  discretized problem definition (identical mesh, first-order
  N\'ed\'elec elements, identical lumped-shunt junction) and are otherwise
  independent.}
  \label{tab:architectures}
  \footnotesize
  \setlength{\tabcolsep}{4pt}
  \begin{tabular}{@{}lll@{}}
    \toprule
     & \textbf{geode-fem} & \textbf{Palace} \\
    \midrule
    Language / substrate & Rust / Burn tensors (f64) & C++ / MFEM \\
    Kernel provenance & general-purpose ML tensor stack &
      domain-specific FEM stack (libCEED-class) \\
    Assembly & batched-tensor N\'ed\'elec (sparse global) &
      MFEM $H(\mathrm{curl})$ \\
    Eigensolver & shift-invert Lanczos & SLEPc-class Krylov \\
    Inner linear solve & faer sparse LU (direct) & AMS-preconditioned iterative \\
    Spurious-mode control & none on headline path (\S\ref{sec:spurious-mode}) &
      divergence-free projection \\
    Parallelism & single process & MPI-distributed \\
    \bottomrule
  \end{tabular}
\end{table}

geode-fem assembles the full-tensor N\'ed\'elec system in double precision
via the batched element kernels of Section~\ref{sec:architecture}, reduces
to interior DOFs, and solves the pencil with a real shift-invert Lanczos
iteration whose inner solves are direct sparse-LU factorizations. It runs
in a single process with no intra-solve parallelism today. Palace
assembles via MFEM, applies a divergence-free projection to confine the
Krylov iteration to the physical subspace, preconditions with
auxiliary-space methods (AMS)
\citep{hiptmair-xu,tzaniovkolevpanayotsvassilevski2018parallel}, and
distributes over MPI ranks. These choices have consequences for both
correctness behavior (Section~\ref{sec:spurious-mode}) and performance
(Section~\ref{sec:perf-cpu}); the point of the benchmark is to expose those
consequences honestly rather than to declare a winner.

\section{Validation Methodology}
\label{sec:methodology}

The benchmark protocol is oracle-first: every claim is anchored either to
an exact analytic quantity or to a committed reference artifact, and the
test suite is designed so that silent pipeline failures are caught by
construction.

\paragraph{Same-mesh, same-physics protocol.}
The mesh is generated once, committed, and pinned by SHA-256. Both solvers
consume it at matched (first) order with the identical junction model
(Section~\ref{sec:junction}). This removes meshing and model entry as
confounders: any residual disagreement is attributable to formulation or
solver implementation.

\paragraph{Analytic tripwires.}
Two exact quantities guard against coincidental or self-consistent-but-wrong
agreement. First, the isolated junction LC resonance
$f_{LC} = 1/(2\pi\sqrt{LC}) = 17.60$\,GHz anchors the expected location of
the junction mode. Second, a Josephson inductance-scaling tripwire: doubling
$L$ must move the junction mode by exactly $1/\sqrt{2}$. In our runs the
junction mode moves from $17.4901$ to $12.37$\,GHz, a ratio of $0.7071$ ---
$1/\sqrt{2}$ to the reported precision (Figure~\ref{fig:lscaling}); the
measured tripwire values are committed alongside the benchmark results
(Section~\ref{sec:repro}).

\paragraph{Inverse tests that must fail.}
Alongside tests asserting agreement, the suite contains perturbed
configurations (such as the doubled-$L$ case) that are \emph{required to
disagree} with the committed oracle. A pipeline that still ``passes'' under
such a perturbation is not measuring anything; the inverse tests make that
failure mode loud.

\paragraph{Participation-based mode identification.}
Modes are identified, not eyeballed --- but the identification must respect
the fact that the two solvers expose \emph{different} participation
diagnostics. Within geode-fem, for each eigenvector $x$ the
stiffness-participation ratio
\begin{equation}
  p \;=\; \frac{x^{\mathsf{T}} K_{\mathrm{port}}\, x}
               {x^{\mathsf{T}} \left(K + K_{\mathrm{port}}\right) x}
  \label{eq:participation}
\end{equation}
measures the fraction of pencil energy stored in the junction inductor and
cleanly separates the junction LC mode ($p = 1.000$) from the resonator and
box modes ($p = 0.000$ at the artifact's three-decimal precision). Palace
reports a per-port energy-participation (port-EPR) output, which is a
field-based, differently normalized quantity; on the committed run the two
diagnostics do \emph{not} rank the modes the same way (in Palace's
committed \texttt{port-EPR.csv} the $17.49$\,GHz junction mode carries the
smallest-magnitude port-EPR of the six modes, all of which are
$\leq 5\times10^{-4}$). The two quantities are complementary diagnostics,
not a shared discriminant, and we do not compare them numerically.
Cross-solver mode identification therefore rests on \emph{frequency}
agreement --- the junction mode appears at $17.490$\,GHz (geode-fem) and
$17.4901$\,GHz (Palace), matching the analytic anchor
$f_{LC} = 17.60$\,GHz --- with geode-fem's participation used to identify
the junction mode on its own side and to expose the spurious mode of
Section~\ref{sec:spurious-mode}.

\paragraph{Determinism check.}
The Palace oracle is not a one-off: rerunning Palace from the committed
configuration reproduced its committed \texttt{eig.csv} bit-for-bit.

\section{Results: Cross-Solver Agreement}
\label{sec:agreement}

Table~\ref{tab:agreement} presents the headline result: on the identical
133k-DOF mesh at matched order, all six modes agree to $\leq 0.033\%$
(worst case $0.032\%$), and the junction LC mode --- the mode the junction
model exists to produce --- agrees to $0.001\%$. This is the evidence
standard the architecture claim of Section~\ref{sec:intro} is held to: a
solver whose assembly and operators are batched tensor algebra on an ML
stack reproduces the domain-specific reference implementation at
production accuracy on a real device geometry.

\begin{table}[t]
  \centering
  \caption{Eigenmode agreement on the identical mesh (22{,}684 nodes /
  133{,}314 tets / 156{,}863 N\'ed\'elec DOFs, 133{,}108 interior after
  PEC), first order in both solvers, junction as a lumped reactive shunt
  ($L = 14.860$\,nH, $C = 5.5$\,fF). geode-fem at commit
  \texttt{3174015}; Palace at commit \texttt{fba6a5b}. All values are
  quoted verbatim from the committed
  \texttt{benchmarks/transmon\_eigen/results.toml} (mode keys in the first
  column): geode-fem frequencies at the committed three decimals, Palace at
  full precision, and the committed per-mode \texttt{rel\_err\_pct}, which
  is computed against the rounded geode-fem values (precision disclosure in
  Section~\ref{sec:repro}).}
  \label{tab:agreement}
  \small
  \begin{tabular}{@{}llccc@{}}
    \toprule
    Mode & & geode-fem (GHz) & Palace (GHz) & $\Delta$ (\%) \\
    \midrule
    \texttt{resonator} & readout resonator & 5.153 & 5.151335830348 & 0.032 \\
    \texttt{mode\_2} & box/cavity & 15.465 & 15.46052107794 & 0.029 \\
    \texttt{junction\_lc} & junction LC & 17.490 & 17.49010903536 & 0.001 \\
    \texttt{mode\_4} & box/cavity & 18.693 & 18.69165792915 & 0.007 \\
    \texttt{mode\_5} & box/cavity & 20.703 & 20.69755679425 & 0.026 \\
    \texttt{mode\_6} & box/cavity & 26.088 & 26.08089940472 & 0.027 \\
    \bottomrule
  \end{tabular}
\end{table}

\begin{figure}[t]
  \centering
  \anvilfig{fig2-agreement.pdf}{fig2\_agreement.py}
  \caption{Mode-frequency agreement: geode-fem versus Palace on the
  identical mesh, with per-mode relative deviations annotated. Data from
  the committed benchmark artifacts
  (\texttt{benchmarks/transmon\_eigen/results.toml} and Palace's committed
  \texttt{eig.csv}).}
  \label{fig:agreement}
\end{figure}

\begin{figure}[t]
  \centering
  \anvilfig{fig3-lscaling.pdf}{fig3\_l\_scaling.py}
  \caption{Josephson inductance-scaling tripwire: junction-mode frequency
  versus junction inductance $L$ on log-log axes. Doubling $L$ moves the
  junction mode from $17.4901$ to $12.37$\,GHz --- ratio $0.7071 =
  1/\sqrt{2}$ --- exactly as the $f \propto 1/\sqrt{L}$ scaling requires.
  Measured values from the committed
  \texttt{benchmarks/transmon\_eigen/results.toml} tripwire block.}
  \label{fig:lscaling}
\end{figure}

Three features of Table~\ref{tab:agreement} deserve comment. First, the
junction mode sits at $17.490$\,GHz (geode-fem) and $17.4901$\,GHz
(Palace), $0.6\%$ below the isolated-circuit anchor
$f_{LC} = 17.60$\,GHz --- the small shift reflects the junction's embedding
in the field problem, and its near-coincidence with the analytic value in
\emph{both} solvers is the tripwire working as intended. Second, the mode
identification is robust: within geode-fem, junction participation is
$p = 1.000$ for the $17.49$\,GHz mode and $p = 0.000$ for every other
physical mode, and the cross-solver assignment rests on frequency agreement
against the analytic anchor rather than on comparing the two solvers'
differently normalized participation outputs
(Section~\ref{sec:methodology}). Third, the agreement is between genuinely
independent code paths --- a batched-tensor assembly feeding a
direct-factorization Lanczos in Rust, and a projected, preconditioned
Krylov method in C++ --- so a common implementation error is an unlikely
explanation for the residual $\leq 0.033\%$ (worst case $0.032\%$), and
the analytic tripwires exclude self-consistent-but-wrong agreement about
the junction physics.

\section{The Absent Qubit Mode: An Expected Negative}
\label{sec:absent-mode}

A reader familiar with the DeviceLayout.jl blog workflow
\citep{devicelayout-blog} will notice that its reported spectrum spans
$[4.14,\ 5.591]$\,GHz --- including a ${\sim}4$\,GHz qubit mode --- while
Table~\ref{tab:agreement} contains no mode below $5.15$\,GHz. This is not a
discrepancy to be minimized; it is a property of the model, and we state it
plainly: \textbf{no ${\sim}4$\,GHz qubit mode exists in this junction model,
by construction.}

The physical transmon frequency arises from the Josephson inductance
resonating against the \emph{pad (shunt) capacitance} of roughly
$80$--$100$\,fF \citep{koch2007}. The v1 lumped-shunt model of
Section~\ref{sec:junction} gives the junction only its own $5.5$\,fF, so
the junction's resonance lands at $17.5$\,GHz, not near $4$\,GHz; the pad
capacitance exists in the field solution but is not wired across the
lumped element. Both solvers agree on the absence: Palace's projected
eigensolve, instructed to hunt at a shift of $\sigma = 4.5$\,GHz, finds
nothing below the $5.15$\,GHz resonator mode. The blog's
$[4.14,\ 5.591]$\,GHz spectrum is therefore \emph{not reproduced by this
model}, and no agreement with it is claimed.

The route to genuine qubit quantities from this benchmark is
energy-participation-ratio (EPR) post-processing over the field modes
\citep{minev2021,nigg2012}, in which the junction's participation in each
electromagnetic mode --- exactly the quantity in
Eq.~\ref{eq:participation} --- feeds the quantization. A first quantum
layer along that route has since been built and measured on this same
mesh; Section~\ref{sec:discussion} summarizes it in one forward-looking
paragraph, and the full analysis is out of scope for the present
cross-validation.

\section{A Disclosed Spurious Mode: Diagnosed, Characterized, and Resolved}
\label{sec:spurious-mode}

Honest cross-validation requires reporting where the solvers
\emph{disagree}, and there is exactly one such place. geode-fem's headline
Lanczos iteration operates on the raw pencil of Eq.~\ref{eq:pencil}
without a gauge projection. On this problem it leaks one spurious mode at
$3.4528$\,GHz, below the physical band, localized on the junction surface
with participation $p = 0.9942$. Palace, whose eigensolve projects onto
the divergence-free subspace, shows no such mode --- its spectrum is clean
down through the $\sigma = 4.5$\,GHz hunt of
Section~\ref{sec:absent-mode}. The mode is excluded from
Table~\ref{tab:agreement} by documented criteria (absence from the
projected Palace spectrum on the identical mesh; location far below the
analytic junction anchor while carrying near-total junction
participation). In an earlier version of this work that disclosure-plus-
filtering was where the story ended; since then the artifact has been
run to ground through a measured three-step gauge/projection arc, every
step committed alongside the benchmark results
(Section~\ref{sec:repro}), and the mode is now diagnosed, characterized,
and --- at the solver level --- resolved.

\paragraph{Step 1: tree--cotree elimination is not spectrum-preserving.}
The classical gauge for the curl--curl \emph{source} problem is
tree--cotree DOF elimination
\citep{albanese1988solution,manges1995generalized}: a spanning-tree
decomposition of the edge graph identifies exactly
$\mathrm{rank}(d^0_{\mathrm{interior}}) = 13{,}747$ gradient DOFs on this
mesh (the count is proven bit-for-bit against the de~Rham gradient rank),
and eliminating them removes both the near-zero gradient cluster and the
$3.4528$\,GHz mode. But for the \emph{generalized eigenproblem} the
physical eigenvectors have nonzero tree-edge components, so the constraint
shifts the spectrum: the resonator moves from $5.1528$ to $5.2372$\,GHz
--- a $1.64\%$ shift, a converged eigenvalue stable across Krylov sizes,
outside our $\leq 1\%$ bar. This honest negative is pinned as a regression
test.

\paragraph{Step 2: bulk divergence-free projection preserves the cavity
but deflates the junction.}
The spectrum-preserving alternative is the M-orthogonal projector
$P = I - G(G^{\mathsf{T}}MG)^{-1}G^{\mathsf{T}}M$ with
$G = d^0_{\mathrm{interior}}$, applied inside the Lanczos iteration ---
structurally the same divergence-free confinement Palace applies
\citep{palace}. Measured on the 133k-DOF fixture, it does what tree--cotree
elimination could not: the cavity spectrum is preserved (resonator at
$0.029\%$ versus Palace, against the gauge's $1.64\%$ shift), every Ritz
vector is divergence-free to ${\sim}10^{-15}$, and the 13{,}747-mode
near-zero gradient cluster collapses to a single survivor. But two
measured facts show the \emph{bulk} projection is incompatible with the
lumped-inductor formulation. First, the spurious $3.4528$\,GHz mode
\emph{survives}: its eigenvector is genuinely solenoidal
(divergence ratio ${\approx}6\times 10^{-15}$), so it is not a
bulk-gradient artifact at all. Second, the \emph{physical} junction LC
mode is deflated away: measured directly on the raw ungauged junction
eigenvector, its divergence ratio
$\lVert G^{\mathsf{T}}Mx\rVert/\lVert x\rVert_M = 50.2$ is macroscopic and
its projected norm $\lVert Px\rVert_M/\lVert x\rVert_M = 1.06\times10^{-4}$
is near zero --- a lumped inductor is a quasi-static, curl-free flux path,
so the junction mode lives almost entirely in the gradient subspace and
the bulk projection cannot distinguish it from the artifacts it removes.

\paragraph{Step 3: port-aware rank-1 re-admission retains all six modes.}
The resolution re-admits exactly one direction: with
$\hat{u} = (I-P)\,x_{\mathrm{junction}} /
\lVert (I-P)\,x_{\mathrm{junction}}\rVert_M$ extracted from the ungauged
junction eigenvector, the rank-1 M-orthogonal update
$P' = P + \hat{u}\hat{u}^{\mathsf{T}}M$ is a genuine projector onto
(divergence-free subspace) $\oplus\ \mathrm{span}\{\hat{u}\}$ --- identity
on the cavity modes \emph{and} on the junction flux, annihilator of the
other 13{,}746 gradient directions. The measured acceptance: all six
physical modes agree with Palace to $\leq 0.029\%$ worst case,
\emph{including} the junction LC mode retained at $17.4901$\,GHz, with the
gradient cluster gone. The headline comparison of
Table~\ref{tab:agreement} deliberately remains the committed ungauged
benchmark (the port-aware solver is a composite --- an ungauged junction
extract feeding the projected band solve --- and shipped after the
headline measurement); its acceptance numbers are committed alongside.

\paragraph{What the spurious mode actually is.}
Two discriminating measurements characterize the $3.4528$\,GHz artifact.
Under the same L-doubling harness as the physical tripwire, it moves to
$2.4449$\,GHz --- ratio $0.7081$, essentially the junction $1/\sqrt{L}$ law
($1/\sqrt{2} = 0.7071$; a box or mesh mode would not move at all) --- and
$99.4\%$ of its stiffness energy sits in the $K_{\mathrm{port}}$ surface
term ($p = 0.9942$). It is therefore a second LC-like resonance of the
$S_\Gamma$ port-surface operator: a $K_{\mathrm{port}}$-driven,
junction-localized eigenmode with no cavity counterpart. Palace's
\texttt{LumpedPort} formulation represents port unknowns differently
(eliminated, rather than added to the full-space stiffness as our
$S_\Gamma$ term), which is exactly why it has no Palace analogue --- and,
being solenoidal, no divergence-free projection, bulk or port-aware,
removes it. On the headline path it remains filtered by the documented
frequency-matching criteria; the difference is that the filter now
excludes a \emph{characterized} artifact rather than an unexplained one.
Figure~\ref{fig:participation} contrasts the two solvers' spectra,
spurious mode included. We regard the arc itself as a finding: it measures
precisely what Palace's divergence-free projection buys in a production
eigensolver, and what a lumped-element formulation demands beyond it.

\begin{figure}[t]
  \centering
  \anvilfig{fig6-participation.pdf}{fig6\_participation.py}
  \caption{geode-fem junction-participation spectrum on the identical mesh,
  with Palace's mode locations (from the committed \texttt{eig.csv})
  overlaid. The physical junction LC mode appears at $17.49$\,GHz in both
  solvers' spectra (geode-fem stiffness participation $p = 1.000$);
  geode-fem's un-projected headline Lanczos additionally leaks one spurious
  junction-localized mode at $3.4528$\,GHz ($p = 0.9942$) with no Palace
  counterpart --- characterized in Section~\ref{sec:spurious-mode} as a
  port-surface artifact obeying the junction $1/\sqrt{L}$ law. Palace's
  field-based port-EPR is a differently normalized diagnostic and is not
  plotted on geode-fem's participation axis
  (Section~\ref{sec:methodology}).}
  \label{fig:participation}
\end{figure}

\section{Performance: CPU Cell}
\label{sec:perf-cpu}

\subsection{Setup}

All CPU measurements ran on a single Amazon EC2 \texttt{m6i.4xlarge}
instance (8 physical cores / 16 vCPUs, 64\,GB RAM) in
\texttt{us-west-2}.\footnote{On-demand cost approximately \$0.77/hour at
the time of the runs; the entire CPU cell is reproducible for a few
dollars.} Each measurement covers the \emph{full pipeline} --- mesh load,
assembly, solve, and output --- timed with \texttt{/usr/bin/time};
entries with uncertainties are the mean over $n = 3$ runs. geode-fem ran
at commit \texttt{3174015}, Palace at commit \texttt{fba6a5b}, both
release builds. The complete cell --- per-run wall clocks, peak
resident-set sizes, and instance/software provenance --- is committed as
\texttt{benchmarks/transmon\_bench\_cpu/results.toml}.

\subsection{Results}

\begin{table}[t]
  \centering
  \caption{CPU-cell wall-clock and memory, full pipeline (mesh load +
  assembly + solve + output), \texttt{m6i.4xlarge}. Uncertainties are over
  $n = 3$ runs where shown; the 4-rank Palace row is a single run
  ($n = 1$). Palace memory is the per-rank maximum single-process RSS, not
  an aggregate across ranks. Values from the committed
  \texttt{benchmarks/transmon\_bench\_cpu/results.toml}. Palace at 16
  ranks (hyperthreads) was refused by MPI process binding and is excluded
  --- it is not a data point.}
  \label{tab:cpu}
  \small
  \begin{tabular}{@{}llcc@{}}
    \toprule
    Solver & Parallelism & Wall (s) & Peak RSS \\
    \midrule
    geode-fem \texttt{@3174015} & 1 process (serial faer sparse-LU
      shift-invert Lanczos) & $51.2 \pm 0.4$ & 3.1\,GB \\
    Palace \texttt{@fba6a5b} & 4 MPI ranks & 50.84 & 0.69\,GB/rank \\
    Palace \texttt{@fba6a5b} & 8 MPI ranks & $30.6 \pm 0.1$ &
      ${\sim}0.5$\,GB/rank \\
    \bottomrule
  \end{tabular}
\end{table}

\begin{figure}[t]
  \centering
  \anvilfig{fig4-cpu-wallclock.pdf}{fig4\_cpu\_wallclock.py}
  \caption{CPU-cell wall-clock comparison (geode-fem single process versus
  Palace at 4 and 8 MPI ranks), with a per-core-efficiency inset. Full
  pipeline timings from Table~\ref{tab:cpu}.}
  \label{fig:cpu}
\end{figure}

Table~\ref{tab:cpu} supports two statements that must be read together.
Per core, geode-fem's serial direct-factorization eigensolve is roughly
$4\times$ more efficient than Palace's distributed Krylov--Schur/AMS solve
on this 133k-DOF problem: one geode-fem core matches four Palace ranks
almost exactly ($51.2$ versus $50.84$\,s). At the whole-box level, Palace's
MPI parallelism wins $1.7\times$ ($30.6$ versus $51.2$\,s), because
geode-fem has no intra-solve parallelism today and leaves seven cores
idle.\footnote{The committed cell measures commit \texttt{3174015} and is
not silently updated. Subsequently, a merged exact-arithmetic reordering
of the Lanczos reorthogonalization (caching $M v_j$ to eliminate an
$O(k^2)$ sparse-matvec pattern; repository PR \#510) reduced the same
133k-DOF eigensolve's release-test wall time from $34.9$ to $21.3$\,s
($1.64\times$) with a bit-for-bit identical spectrum. The
\texttt{m6i.4xlarge} full-pipeline cell has not been re-measured
post-\#510, so Table~\ref{tab:cpu} stands as committed; the follow-up
number is reported here as the separate merged fact it is.}

This is an architecture trade-off, not a ranking. At 133k DOFs the sparse
LU factorization fits comfortably in memory (3.1\,GB peak) and its flop
efficiency dominates; a distributed iterative method carries communication
and preconditioner-setup overheads that cannot amortize at this size, but
distributes memory (${\sim}0.5$\,GB/rank at 8 ranks) and scales to problem
sizes and rank counts where a serial direct factorization is not viable at
all. The honest summary is that each architecture wins the regime it was
designed for, and this benchmark documents where the crossover sits for a
production transmon geometry at first order. Palace at 16 ranks on the
hyperthreaded vCPUs was refused by MPI process binding; we exclude it
rather than report a degraded configuration as if it were a measurement.

\section{Performance: GPU Cell (Correctness Established; Scaling an
Honest Negative)}
\label{sec:perf-gpu}

\subsection{Scope}

geode-fem's GPU path today accelerates the \emph{driven} (deterministic
frequency-sweep) solve --- the matrix-free apply and on-device COCG of
Section~\ref{sec:architecture} --- and \emph{not} the eigensolve that
produced the headline comparison of Section~\ref{sec:agreement}, which
remains factorization-bound on CPU. A GPU column for the eigenmode
benchmark would therefore be misleading; the GPU cell targets the driven
pipeline, in two stages: correctness on physical hardware, then a
four-configuration scaling study. All GPU execution is CUDA-f32
(Section~\ref{sec:architecture}), so accuracy against the f64 reference is
reported alongside every timing.

\subsection{Correctness on physical hardware}

On 2026-07-14, geode-fem code executed on a physical GPU for the first
time: an Amazon EC2 \texttt{g6e.xlarge} instance (1$\times$ NVIDIA L40S,
46\,GB; driver 595.71.05, CUDA 13.2). Two CUDA-f32 correctness smokes pass
on this hardware: the matrix-free N\'ed\'elec matvec smoke
(\texttt{matrix\_free\_cuda\_f32\_smoke}, 15.0\,s wall including GPU
initialization and kernel JIT compilation) and the on-device COCG smoke
(\texttt{cocg\_burn\_cuda\_f32\_smoke}, 3.3\,s), both asserting agreement
with the f64 CPU reference within f32-appropriate tolerances. The quoted
wall times are smoke-test durations dominated by startup cost ---
correctness evidence, not performance measurements. At the time of those
smokes the assembled end-to-end \texttt{IterativeMatrixFree} path had no
GPU runtime test (a disclosed gap, repository issue \#499); the scaling
study below closes it, running that full path on the GPU across all
measured sizes.

\subsection{Driven-solve scaling: the measured answer}

The scaling cell (repository issue \#501; committed as
\texttt{benchmarks/gpu\_driven\_scaling/results.toml}) solves one physical
problem --- a $\sigma$-lossy parallel-plate cube with a single lumped
port, from the committed driven-solve fixture family --- at a fixed
frequency across four solver configurations while the mesh is refined
from $1{,}854$ to $25{,}695$ edges. All sixteen cells were measured on the
\emph{same host} (the \texttt{g6e.xlarge} above; 4 vCPUs), so the
GPU-versus-CPU comparison is same-host by construction --- and, for the
same reason, the CPU baselines here are not comparable to the
\texttt{m6i} numbers of Section~\ref{sec:perf-cpu}. This cell is a
driven-solve scaling study on a cube fixture family, not the transmon
eigenproblem. Table~\ref{tab:gpu} reports the committed medians.

\begin{table}[t]
  \centering
  \caption{Driven-solve wall-clock scaling, solve-only medians of $n = 3$
  timed repetitions (warm-up excluded), same host (\texttt{g6e.xlarge}:
  4 vCPU + 1$\times$ L40S). One fixture ($\sigma$-lossy lumped-port cube)
  at fixed $\omega$, mesh-refined across the columns. f64 iterative
  configurations run at tolerance $10^{-8}$; the f32 GPU configuration at
  $10^{-6}$ (its recurrence floor sits above $10^{-8}$), so iteration
  counts are not exactly comparable. The accuracy row is the full-field
  relative $L_2$ error of the GPU-f32 solution against the Direct-f64
  reference at the same size. At 25{,}695 edges the GPU-f32 configuration
  converged in the committed single-frequency runs but its five-point
  frequency sweep did not, and run-to-run convergence is nondeterministic
  at that size (f32 floor). Values verbatim from the committed
  \texttt{benchmarks/gpu\_driven\_scaling/results.toml}.}
  \label{tab:gpu}
  \small
  \setlength{\tabcolsep}{5pt}
  \begin{tabular}{@{}lcccc@{}}
    \toprule
    & \multicolumn{4}{c}{Mesh size (N\'ed\'elec edges)} \\
    \cmidrule(l){2-5}
    Configuration & 1{,}854 & 5{,}859 & 13{,}428 & 25{,}695 \\
    \midrule
    Direct faer LU (CPU f64), s & 0.024 & 0.203 & 1.540 & 6.036 \\
    Assembled-CSR COCG (CPU f64), s & 0.032 & 0.198 & 0.709 & 1.865 \\
    Matrix-free COCG (CPU f64), s & 1.652 & 8.885 & 30.234 & 82.056 \\
    Matrix-free COCG (GPU f32), s & 4.388 & 13.351 & 34.381 & 81.764 \\
    \midrule
    GPU-f32 accuracy vs.\ Direct-f64 (rel.\ $L_2$) &
      $1.26\times10^{-4}$ & $2.93\times10^{-4}$ &
      $4.60\times10^{-4}$ & $1.20\times10^{-3}$ \\
    \bottomrule
  \end{tabular}
\end{table}

The honest read: \textbf{GPU-f32 matrix-free loses to every CPU
configuration at every measured size on this host.} The assembled-CSR COCG
is $136\times$ faster at $1{,}854$ edges ($0.032$ versus $4.39$\,s) and
still $44\times$ faster at $25{,}695$ edges ($1.86$ versus $81.8$\,s);
even the direct LU is $13\times$ faster at the top size ($6.04$\,s). The
mechanism is unambiguous: the GPU's per-iteration cost is kernel-launch
dominated at these sizes --- ${\sim}28$\,ms per iteration at the top size,
against ${\sim}0.63$\,ms for the assembled-CSR path on CPU --- because
$25.7$k edges is far too small to fill an L40S. Accuracy compounds the
verdict: the f32 cells degrade from $1.26\times10^{-4}$ to
$1.20\times10^{-3}$ relative error with refinement (the true recomputed
residual floors at $6.2\times10^{-4}$ to $5.4\times10^{-3}$, orders above
the requested tolerance), and at the top size convergence is
nondeterministic run-to-run --- usable for ${\sim}0.1\%$-class field
quantities, but trending the wrong way for refinement studies.

One directional positive survives, and it is the one the architecture
predicts: against the \emph{same matrix-free algorithm on CPU} (the
identical Burn source path on the ndarray backend), the GPU's deficit
falls monotonically from $2.7\times$ slower at $1{,}854$ edges to parity
at $25{,}695$ ($81.8$ versus $82.1$\,s) --- the device-resident iteration
catches the CPU at ${\sim}26$k edges and would pass it beyond. That
crossover is currently moot in f32, because the top size already sits at
the f32 convergence ceiling. We therefore state the performance position
as trajectory, not achievement, with three concrete levers, each tracked:
(i) kernel fusion and launch batching to attack the ${\sim}28$\,ms
per-iteration launch overhead, the dominant cost at benchmark sizes;
(ii) f64 (or mixed-precision/iterative-refinement) on device when the
cubecl runtime ships it, which removes the convergence ceiling that makes
larger problems unreachable --- noting honestly that the L40S's
$1/64$-rate f64 would trade ceiling for raw rate, and datacenter-class
f64 parts change that arithmetic; and (iii) preconditioning the on-device
COCG to cut iteration counts, which multiplies whatever the first two
levers yield. Nothing in Table~\ref{tab:gpu} contradicts the
device-resident-Krylov design intent; it bounds where that design pays
off --- well above $26$k edges, in f64 --- and we prefer publishing the
bound to implying the win.

% Figure plan slot 5 (GPU cell) intentionally remains a table, not a
% figure: the four-config scaling data is fully carried by Table
% \ref{tab:gpu} and the BRIEF directs no new figure be fabricated for it.

\section{Reproducibility}
\label{sec:repro}

Every artifact behind every number in this paper is committed and
scripted.\footnote{Repository URL and archival DOI for the geode-fem
source at commit \texttt{3174015} to be added at submission.
TODO(operator).}

\begin{itemize}
  \item \textbf{Geometry and mesh.} Generated by DeviceLayout.jl v1.15.0's
    \texttt{SingleTransmon} example \citep{devicelayout}; the fixture
    generator is committed, and the MSH~4.1 fixture is pinned by SHA-256
    (\texttt{5b3ff4c3\ldots{}b33dd}; full hash and provenance in the
    committed \texttt{transmon\_smoke.provenance.txt}).
  \item \textbf{The MSH conversion gotcha.} MFEM's MSH~4.1 reader rejects
    the DeviceLayout multi-entity-block file (``vertices indices are not
    unique''), so the Palace runs consume an MSH~2.2 conversion
    (\texttt{gmsh -save -format msh2}). The conversion is node-preserving
    (all 22{,}684 nodes unchanged, physical groups preserved) and
    physics-neutral: eigenvalues reproduce bit-for-bit across the
    conversion. The conversion step is documented in the committed
    \texttt{palace\_config.provenance.txt}. We document this because it is
    exactly the kind of silent pipeline step a same-mesh claim must
    disclose.
  \item \textbf{Solver configurations.} The Palace configuration and its
    provenance file are committed under
    \path{reference/fixtures/transmon_palace/}; geode-fem's benchmark
    definitions live in \path{benchmarks/transmon_eigen/results.toml},
    alongside the exact benchmark commands.
  \item \textbf{Precision of the committed agreement numbers.}
    \path{benchmarks/transmon_eigen/results.toml} stores geode-fem
    frequencies at three decimals while the Palace column carries full
    precision, so the committed per-mode \texttt{rel\_err\_pct} is
    computed against the rounded geode-fem values. At full precision the
    agreement is slightly better (approximately $0.029\%$ worst case). We
    quote the committed-artifact numbers throughout and disclose the
    rounding here.
  \item \textbf{Performance artifacts.} The complete CPU cell
    (per-run wall clocks, peak RSS, instance and software provenance) is
    committed as \path{benchmarks/transmon_bench_cpu/results.toml}. The
    complete GPU driven-solve scaling cell (all sixteen configuration
    $\times$ size cells, iteration counts, recomputed true residuals,
    accuracy-versus-Direct columns, DNF/flaky handling, and same-host
    provenance) is committed as
    \path{benchmarks/gpu_driven_scaling/results.toml}.
  \item \textbf{Tripwire, gauge-arc, and spurious-mode artifacts.} The
    measured L-doubling tripwire values ($17.4901 \to 12.37$\,GHz, ratio
    $0.7071$), the spurious-mode measurements ($3.4528$\,GHz,
    $p = 0.9942$; L-doubling $3.4528 \to 2.4449$\,GHz, ratio $0.7081$),
    and the full gauge/projection arc of
    Section~\ref{sec:spurious-mode} --- tree--cotree counts and the
    $1.64\%$ shift, the bulk-projection cavity/junction measurements
    (divergence ratio $50.2$; projected norm $1.06\times10^{-4}$), and
    the port-aware acceptance (all six modes, worst case $0.029\%$) ---
    are committed in \path{benchmarks/transmon_eigen/results.toml}, each
    pinned by a release-tier regression test.
  \item \textbf{Oracles.} Palace's raw outputs (\texttt{eig.csv},
    per-port participation, run log) are committed; a rerun from the
    committed configuration reproduced \texttt{eig.csv} bit-for-bit.
  \item \textbf{Hardware disclosure.} CPU cell: EC2 \texttt{m6i.4xlarge}
    (8 physical cores / 16 vCPU, 64\,GB), \texttt{us-west-2}. GPU cell
    (correctness and scaling): EC2 \texttt{g6e.xlarge} (4 vCPU, 30\,GB;
    1$\times$ NVIDIA L40S 46\,GB; driver 595.71.05, CUDA 13.2).
  \item \textbf{Versions.} geode-fem commit \texttt{3174015} (headline
    cells; post-measurement merges are cited by PR number where
    discussed); Palace commit \texttt{fba6a5b}; DeviceLayout.jl v1.15.0;
    Burn/burn-cuda 0.21.
\end{itemize}

\section{Discussion}
\label{sec:discussion}

\paragraph{What the cross-validation does and does not establish.}
In V\&V terms \citep{roache1997quantification,oberkampf2010verification},
this study is code verification and code-to-code comparison, not
validation against nature: two independent implementations agreeing to
$\leq 0.033\%$ (worst case $0.032\%$) on the identical discrete problem,
with analytic tripwires pinning the junction physics. It does not certify
that either solver's answer matches a measured device --- that is what
\citet{ye2025electromagnetic}-style experimental benchmarks address, and
the two evidence types are complementary. What same-mesh agreement
uniquely rules out is the large class of errors that live in formulation,
assembly, boundary-condition handling, and eigensolver machinery --- and
for the architecture thesis specifically, it rules out the possibility
that phrasing assembly and operators as batched ML-stack tensor algebra
silently costs accuracy on a production geometry.

\paragraph{What the performance evidence supports.}
The measured position is: per-core CPU competitiveness demonstrated
(${\sim}4\times$ versus the distributed reference at 133k DOFs), GPU
portability demonstrated at essentially zero incremental engineering (the
CUDA leg is the same source path as the CPU leg, a backend type parameter
away), and GPU \emph{advantage} not demonstrated --- the scaling cell of
Section~\ref{sec:perf-gpu} is an honest negative at benchmark sizes, with
a measured parity crossover against the same algorithm on CPU at
${\sim}26$k edges and named, tracked levers beyond it. The claim this
paper stands behind is therefore viability-plus-trajectory, not achieved
GPU speedup.

\paragraph{Threats to validity.}
A common-mode error --- both solvers wrong in the same way --- cannot be
excluded by cross-comparison alone; the analytic anchors ($f_{LC}$, the
$1/\sqrt{2}$ scaling) and the genuinely disjoint implementations make it
unlikely for the junction physics specifically. Configuration-entry
mismatch (comparing subtly different problems) is mitigated by committed
configurations, the shared mesh file, and participation-based mode
identification on the geode-fem side with frequency-anchored cross-solver
assignment (Section~\ref{sec:methodology}). Performance numbers are $n = 3$
measurements on one instance type per cell; they support the
architecture-level readings in
Sections~\ref{sec:perf-cpu}--\ref{sec:perf-gpu}, not microsecond-level
claims, and the GPU cell's CPU baselines are same-host (4-vCPU) values
that must not be compared across boxes.

\paragraph{Toward qubit quantities.}
A first quantum layer above this benchmark has since been built and
measured on the same mesh (committed as
\path{benchmarks/transmon_quantum/results.toml}): a tensor-$\varepsilon$
electrostatic solve on the mesh's three node-disjoint metal conductors
yields the island Maxwell capacitance $C_\Sigma = 136.7$\,fF, hence
$E_C = 0.142$\,GHz and $E_J/E_C = 77.6$, and the charge-basis-exact
transmon spectrum of \citet{koch2007} gives $\omega_{01} = 3.38$\,GHz and
anharmonicity $\alpha = -0.158$\,GHz --- a plausible transmon, computed
end-to-end from the same committed geometry. It carries its own honest
finding: $C_\Sigma$ exceeds the ${\sim}90$\,fF design-anchor value, and
the discrepancy is diagnosed as a model-definition difference rather than
a truncation artifact (grounding the exterior wall moves $C_\Sigma$ by a
relative $6\times10^{-5}$; the scalar-versus-tensor sapphire treatment
shifts it $0.75\%$). Connecting this electrostatic route to the
field-mode EPR quantization \citep{minev2021,nigg2012} over the
eigenmodes of Section~\ref{sec:agreement} is a separate paper's worth of
analysis and is deliberately not compressed into this one.

\paragraph{Limitations.}
(i) One geometry, one mesh, first order only --- higher-order and
mesh-refinement convergence is future work (a $p = 2$ Palace
configuration is already committed for that purpose). (ii) The lossless v1
model leaves readout ports open; resistive terminations (and hence
complex-symmetric pencils and mode $Q$s) are out of scope. (iii) The
headline eigenpath remains the ungauged committed benchmark; the
port-aware projection of Section~\ref{sec:spurious-mode} retains all six
modes at $\leq 0.029\%$ but is a composite solver (an ungauged extract
feeds it), and the solenoidal port artifact is characterized and filtered
rather than eliminated. (iv) The per-core efficiency result is specific
to the 133k-DOF regime where a serial direct factorization is
memory-feasible; it will not extrapolate to problem sizes that demand
distribution. (v) The GPU evidence covers driven-solve correctness and an
honest scaling negative on a cube fixture family well below
GPU-saturating sizes, all in f32; there is no GPU eigensolve path, and no
GPU claim in this paper extends beyond what
Table~\ref{tab:gpu} contains.

\section{Conclusion}
\label{sec:conclusion}

We set out to test whether a general-purpose ML tensor-compiler stack can
carry production-grade computational electromagnetics, and to hold the
answer to a cross-validation standard rather than a demonstration
standard. The evidence: a Burn-native full-wave $H(\mathrm{curl})$ solver
whose assembly is batched tensor algebra reproduces Palace --- the
domain-specific reference stack, itself built on a FEM-community element
JIT --- to $\leq 0.033\%$ across all six eigenmodes of a real transmon
geometry on the identical mesh (worst case $0.032\%$; $0.001\%$ on the
junction LC mode), with exact analytic tripwires guarding the junction
physics; it matches the reference on serial efficiency (${\sim}4\times$
per core at 133k DOFs, against a $1.7\times$ whole-box win for MPI
distribution); and it gains a portable CUDA execution path from the same
source, whose measured scaling cell we report as the honest negative it
currently is --- kernel-launch-bound losses at every benchmark size, a
parity crossover against its own algorithm on CPU at ${\sim}26$k edges,
and named levers (fusion, on-device f64, preconditioning) between here
and advantage. Along the way the benchmark delivered the first
independent validation of the Palace/DeviceLayout.jl transmon stack, a
reusable oracle-first cross-solver methodology, and an honest-physics
record upgraded from disclosure to diagnosis: the by-construction absent
qubit mode confirmed by both solvers, and the spurious mode run to ground
through a measured gauge/projection arc that ends with all six physical
modes retained at $\leq 0.029\%$ and the artifact characterized as a
port-surface resonance. The architectural conclusion is measured but
real: the ML ecosystem's compiler investment is inheritable by
finite-element electromagnetics at production accuracy today, with the
performance frontier --- not the correctness frontier --- as the open
work.

\paragraph{Acknowledgments.}
geode-fem's development was AI-orchestrated (agent-driven implementation
with human direction and review); this benchmark program and paper follow
the same discipline, with every reported number traced to a committed,
human-auditable artifact.
% TODO(operator): finalize acknowledgment wording before submission.
% TODO(operator): whiteroom L1-L4 operator specification — decide
% public/citable (repo/DOI) vs acknowledged-not-cited (BRIEF coauthor note).

\bibliographystyle{plainnat}
\bibliography{refs}

\end{document}
